% !TeX root = ../Book.Stress_regulation.tex
\documentclass[../main.tex]{subfiles}
\graphicspath{{\subfix{../images/}}}
\begin{document}

\chapter{The Rune Exercises}\index{exercises!Runes}\index{exercises!energy}

\section{Rune Exercise 1: U--Position}

\label{Ex:RuneU}
\begin{figure}[htb!]
	\centering
	\includegraphics[height=10cm]{Rune_U.jpg}\label{ex:runeU}
	\caption{The Rune Exercises: Rune U}
\end{figure}


This bending position exists in many systems, as for instance in yoga or the rune exercises. 
They stimulate the head--nerve center.

\subsection*{Starting Position}

Stand with your feet parallel and shoulder width apart. 
Then slowly lift your arms above your head.

\subsection*{Execution}

Inhale as deeply as possible and then slowly exhale. 
Bend simultaneously the torso and head forward and down, till the torso is in a horizontal position.
The arms and finger should point downwards, towards the earth (see picture~\ref{ex:runeU}).
Slightly bend your knees, in order for your back to be more relaxed.
Shift your weight more and more towards your hands.
Stay in this position for the duration of 10 breaths. Increase this length daily by one breath, till you reached an optimal duration of 30 breaths.

If you are not able to execute the exercise due to health reasons, then execute the exercise while being seated. 
In this case, sit with a straight spine and lift your arms above your head, while calmly and deeply inhaling.
While exhaling, bend your torso as much forward as possible. 
Let your arms hang towards the floor, left and right of your legs.
Breathe and increase the length as described above.

\newpage

\section{Rune Exercise 2: I--Position}

\label{Ex:RuneI}
\begin{figure}[htb!]
	\centering
	\includegraphics[height=10cm]{Rune_I.jpg}\label{ex:runeI}
	\caption{The Rune Exercises: Rune I}
\end{figure}

The exercise of the rune I served the Germanic peoples to increase the mental and physical power of resistance.
Especially his exercise is also known as the power exercise. In measurements it is shown as highly effective at relaxing the main interference fields and muscle tension. 

\subsection*{Starting Position}

Stand up or sit, if possible facing the west.

\subsection*{Execution}

Slowly lift your hands above your head.
The palm of your hands are facing each other, in a distance of about 2/3 of a foot (20 cm, see picture~\ref{Ex:RuneI}).
Breathe slowly and deeply in a regular fashion. 
Stay in this position for 10 breaths. Slowly increase this duration by one breath more per day, till you reach the ideal practice time of 30 breaths.

\newpage

\section{Rune exercise 3: Y--Position}

\label{Ex:RuneY}
\begin{figure}[htb!]
	\centering
	\includegraphics[height=10cm]{Rune_Y.jpg}\label{ex:runeY}
	\caption{The Rune Exercises: Rune Y}
\end{figure}
 
 The Y--exercise comes from the rune exercises. It serves in ancient times to transmutate sexual energy and to unlock higher spheres and the subtle worlds.
 
 \subsection*{Starting Position}
 
Continue from the position of exercise 2, without stopping.
 
 \subsection*{Execution}
 
 Let both arms slowly sink down to the point where your body looks like the letter "Y". 
 Your palms are flat and turned upwards (see picture`\ref{ex:runeY}).
 In this exercise pay attention to breathe in a relaxed and calm way. 
 Stay in this position at first for 10 breaths and slowly increase till length by one breath a day to eventually achieve 30 breaths.
 
 \newpage
 \section{Rune exercise 4: F--Position}
 
 \label{Ex:RuneF}
 \begin{figure}[htb!]
 	\centering
 	\includegraphics[height=10cm]{Rune_F.jpg}\label{ex:runeF}
 	\caption{The Rune Exercises: Rune F}
 \end{figure}
 
 This pointing position comes from  (Hatha--)Yoga and is also found in the rune rites. This position helps the growth of your mental powers and looses blockages and energy congestion.
 
 \subsection*{Starting Position}
 
 Continue from the position of exercise 3 without stopping into the exercise 4.
 
 \subsection*{Execution}
 
 
 \label{Ex:RuneF_Hand}
 \begin{figure}[htb!]
 	\centering
 	\includegraphics[height=3cm]{Rune_F_Hands}\label{ex:runeF_Hands}
 	\caption{The Rune Exercises: Rune F, hand position}
 \end{figure}
 
 Lower both your straight arms forward. 
 The left hand should be 4 inches (10 cm) higher up than the right one.
 
 With the fingers, you build the following shape:
 Put the thumb and middle finger (not the indices!) gently together with the tips, just as much that they slightly touch.
 the other ginger are straight forward and pointing from your body in direction west (see picture~\ref{Ex:RuneF_Hand}).
 
 
 
 Then hold your hands and finger in this position, again for the duration of 10 breaths. 
 As with the other exercises, slowly increase the duration to 30 breaths.
 Afterwards, let both hands sink down and consciously relax your shoulders (eventually by slightly moving and loosening) and count mentally to 20.
 Make sure in this exercise, that your finger don't point towards a person.
 
\newpage
  \section{Rune exercise 5: T--Position}
 
 \label{Ex:RuneT}
 \begin{figure}[htb!]
 	\centering
 	\includegraphics[height=10cm]{Rune_T.jpg}\label{ex:runeT}
 	\caption{The Rune Exercises: Rune T}
 \end{figure}
 
 This exercise comes from the rune rites and serves to focus on the cosmic waves in the gland system and solar plexus. This exercise also brings a deep relaxation of the muscles.
 
 \subsection*{Starting Position}
 
 Continue from the position of exercise 4 (F--position) without stopping into the exercise 5.
 
 \subsection*{Execution}
 
 Make with your hands a T--position, where both of your arms should be straight and diagonally pointing downwards (see picture~\ref{Ex:RuneT}). The palms point toward the floor. 
 Maintain this position for the duration of 10 breaths and daily increase the duration by 1 breath up to 30 breaths.
 
  \newpage
  \section{Rune exercise 6: W--Position}
 
 \label{Ex:RuneW}
 \begin{figure}[htb!]
 	\centering
 	\includegraphics[height=10cm]{Rune_W.jpg}\label{ex:runeW}
 	\caption{The Rune Exercises: Rune W}
 \end{figure}
 
 This exercise is known in the martial art Taekwondo as an exercise which is very efficient at relaxing the shoulder and neck muscles. The starting position is also found in the rune exercises. This exercise increases the flow of energy and helps harmonizing your breath.
 
 \subsection*{Starting Position}
 
 Take both hands in front of your body in the height of your face, the edges of your hands facing forward, away from you.
 Hold your hands straight and thumbs flush with your hand, like a karate fighter.
 the edges of your hands are about a foot apart from each other.
 
 \subsection*{Execution}
 While forcefully exhaling through the nose, snap your right arm and edge of the hand to the side. Your arm takes a horizontal position (see picture~\ref{Ex:RuneW}).
 Simultaneous to a snap of your left hand, take you right hand back into the starting position.
 The same movement gets now executed with the left arm.
 Each movement should be executed en times each. 
 it can be gradually increased to 15 movements with each arm.
 
 
 This concludes the first part of the training. let your hands calmly sink down. Breathe a few times in a calm regular way. Consciously relax yourself, all your muscles and let your arms hang down by their own weight.
 
 On the next page, the rune exercises are again condensed and then on the next page is a summary.
 
 \newpage
 \subfile{Rune_exercises1page}
 
\end{document}